\begin{examplebox}
	\textbf{\textcolor{CExample}{例1（基础直接应用）}}

	\textbf{\textcolor{CExample}{题目：}}
	已知抛物线$y^2=4x$的弦$AB$的中点为$M(2,1)$，求弦$AB$所在直线的方程。

	\textbf{\textcolor{CExample}{解题步骤：}}
	\begin{description}[style=nextline, leftmargin=2.8em, labelsep=0.5em, itemsep=0.45em]
		\item[\textbf{第一步（确定参数）：}]将$y^2=4x$与标准方程$y^2=2px$对比，得$2p=4$，故$p=2$。
			\par\textbf{理由：}为应用公式需明确$p$。
		\item[\textbf{第二步（套用公式）：}]由中点弦斜率公式$k=\dfrac{p}{y_0}$，代入$p=2,\ y_0=1$得
			\[
				k = \frac{2}{1}=2
			\]
			\par\textbf{理由：}直接运用结论。
		\item[\textbf{第三步（点斜式写）：}]利用点斜式，直线过$M(2,1)$，斜率$2$，得
			\[
				y-1 = 2(x-2) \quad\Longrightarrow\quad y=2x-3
			\]
			\par\textbf{理由：}求得方程即弦所在直线。
	\end{description}

	\textbf{\textcolor{CExample}{关键结论：}}
	弦$AB$的方程为$y=2x-3$。

	\medskip
	\textbf{\textcolor{CExample}{例2（稍变形应用）}}

	\textbf{\textcolor{CExample}{题目：}}
	已知抛物线$y^2=8x$的弦$AB$的斜率为$-2$，求弦$AB$中点$M$的纵坐标。

	\textbf{\textcolor{CExample}{解题步骤：}}
	\begin{description}[style=nextline, leftmargin=2.8em, labelsep=0.5em, itemsep=0.45em]
		\item[\textbf{第一步（确定参数）：}]由$y^2=8x$得$2p=8$，即$p=4$。
			\par\textbf{理由：}求$p$以代入公式。
		\item[\textbf{第二步（公式变形）：}]中点弦斜率公式$k=\dfrac{p}{y_0}$可变形为$y_0=\dfrac{p}{k}$（$k\neq0$）。
			\par\textbf{理由：}已知斜率求中点坐标需使用逆推形式。
		\item[\textbf{第三步（代入求值）：}]将$p=4,\ k=-2$代入$y_0=\dfrac{p}{k}$得
			\[
				y_0 = \frac{4}{-2} = -2
			\]
			\par\textbf{理由：}直接计算得中点纵坐标。
	\end{description}

	\textbf{\textcolor{CExample}{关键结论：}}
	弦$AB$的中点纵坐标为$y_0=-2$。

	\medskip
	\textbf{\textcolor{CExample}{例3（综合应用）}}

	\textbf{\textcolor{CExample}{题目：}}
	已知抛物线$y^2=2x$，斜率为$k$（$k\neq0$）的直线$l$过定点$P(1,0)$，与抛物线交于$A$、$B$两点，弦$AB$的中点为$M$。求点$M$的轨迹方程。

	\textbf{\textcolor{CExample}{解题步骤：}}
	\begin{description}[style=nextline, leftmargin=2.8em, labelsep=0.5em, itemsep=0.45em]
		\item[\textbf{第一步（设参表示）：}]设直线$l: y = k(x-1)$，设$M(x_0,y_0)$，则$y_0 = \frac{y_1+y_2}{2}$。
			\par\textbf{理由：}引入参数表示直线和中点。
		\item[\textbf{第二步（套用结论）：}]由抛物线$y^2=2x$知$p=1$，根据中点弦斜率公式得$k = \dfrac{p}{y_0} = \dfrac{1}{y_0}$。
			\par\textbf{理由：}直接应用结论建立$k$与$y_0$的关系。
		\item[\textbf{第三步（消参得轨）：}]因为$M$在直线$l$上，故$y_0 = k(x_0-1)$。代入$k=\dfrac{1}{y_0}$得
			\[
				y_0 = \frac{1}{y_0}(x_0-1) \quad\Longrightarrow\quad y_0^2 = x_0 - 1.
			\]
			当$y_0=0$时对应点$P(1,0)$，但此时弦不存在（直线与抛物线相切或退化为一点），故轨迹中应除去该点。
			\par\textbf{理由：}消去参数$k$，并考虑实际意义。
	\end{description}

	\textbf{\textcolor{CExample}{关键结论：}}
	中点$M$的轨迹方程为$y^2 = x - 1$（$y \neq 0$），即除去点$(1,0)$的抛物线段。
\end{examplebox}
